\documentclass[12pt, titlepage]{article}
%\usepackage[norsk]{babel}
\usepackage[utf8]{inputenc}
\usepackage[T1]{fontenc}
\usepackage{minted, enumerate}
\usepackage[colorlinks = true,
linkcolor = blue,
urlcolor = blue]{hyperref}


\begin{document}

\title{Innlevering: Mappe}
\author{IIRA2001 - Høst 2023}
\date{}
\maketitle

Mappen består av de to mappeoppgavene vi har jobbet med i høst og et refleksjonsnotat

\section*{Del I}

I første del av mappen skulle dere skrive 3 pythonprogrammer:
\begin{itemize}
   \item {Et som bruker den logistiske likningen, feks ved populasjonsvekst}
   \item {En sparekalkulator}
   \item {Og en lånekalkulator}
\end{itemize}
Se egen oppgavebeskrivelse på blackboard for mer detaljer om programmene

~\\
\textit{\textbf{MERK:} Dere ble også bedt om å skrive et refleksjonsnotat i del I. Dette skal ikke leveres.
Dere leverer i stedet \textbf{ett} refleksjonsnotat som gjelder hele semesteret}

\section*{Del II}

I andre del av mappen skulle dere lage 2 programmer som analyserte datasett, brukte et web-api eller begge deler.
I hovedsak ville vi ha 2 programmer som gjorde noe interessant med 1 eller flere datasett, hvorav det ene programmet brukte et web-api.
Se oppgavebeskrivelse på blackboard for detaljer om programmene og unntak til beskrivelsen over

\section*{Beskrivelse av programmene}

\begin{itemize}
   \item{Til hvert pythonprogram skulle dere også å skrive litt tekst. }
   \item {Teksten må inneholde en beskrivelse av hva målet med programmet er, feks hvilken problemstilling du undersøker.}
   \item {I tillegg trenger du en liten diskusjon og/eller beskrivelse av resultatet av programmet. }
\end{itemize}

Kort fortalt trenger vi en liten rapport som beskriver hva målet med koden din var, og diskuterer hva resultatet av koden din ble, eller hva du fant ut.\\
Hvor lang teksten er avhenger av programmet ditt. 
Dersom du har laget en sparekalkulator og bruker den til å finne ut hvor mye du må spare per måned for å kjøpe deg en lamborghini i 30-årsgave til deg selv, er det kanskje færre vurderinger og diskusjoner å skrive enn om du undersøker effekten strømprisen i europa har på aksjekursen til norsk hydro.

~\\
Dere kan  levere kode og «rapport» på 2 måter:
\begin{itemize}
   \item {Som 1 *.ipynb fil (Jupyter-notebook fil) per program med kodeceller for python og tekstceller for beskrivelse/resultat/diskusjon}
   \item {Som 1 *.py fil (pythonscript) med koden din og 1 PDF-fil med beskrivelse/resultat/diskusjon, per program}
\end{itemize}

\color{blue}
\subsubsection*{ENDRING}
Etter å ha sett på løsningene deres og gitt tilbakemeldinger, har jeg merket meg at mange har andre fornuftige måter å organisere mappen sin på.
Dersom dere vil lage ett dokument med refleksjonsnotat,  beskrivelse og diskusjon av programmene går det også fint.
Noen har presantert mappen med streamlit og andre mer avanserte måter - dette er også greit, men vær for all del forsiktig med å sørge for at sensor klarer å kjøre programmet/applikasjon deres.

Dere har frihet til å presentere mappen deres på litt andre måter enn slik det er forskrevet over, men på eget ansvar:
\begin{itemize}
   \item {Vi må kunne kjøre koden deres uten for mye styr}
   \item {Koden bør være skilt ut i egne filer}
   \item {Programmene trenger tekst som beskriver de, og diskuterer eventuelle problemstillinger og resultater}
\end{itemize}

\color{black}


\section*{Refleksjonsnotat}

Det siste som skal leveres er et refleksjonsnotat i en PDF-fil.

~\\
I løpet av semesteret har vi sett på grunnleggende programmering i python, og anvendelse av dette i dataanalyse og økonomiske fag.
Siste del av mappen består av et \textit{refleksjonsnotat} hvor dere reflekterer rundt læringsprosessen dere har vært gjennom, og hvordan programmering kan være faglig relevant i videre studie og virke.

Et refleksjonsnotat er en personlig tekst hvor du reflekterer over egne tanker, erfaringer eller innsikter knyttet til en bestemt hendelse eller prosess (slik som læringsprosessen dere nå har vært gjennom).
Ved å bearbeide og reflektere rundt disse tankene og erfaringene er målet å fremme læring og personlig og faglig vekst.

~\\
I hovedtrekk skal dere reflektere rundt hvordan programmering kan være faglig relevant, hvilke ferdiheterer dere har tilegnet dere og hvordan dere ser for dere fremtidig utvikling og bruk av det dere har lært.\\
%\textbf{Tips, relevante punkter å ha med, og relevant spørsmål å reflektere rundt:}
\subsubsection*{Tips, relevante punkter og spørsmål}
\begin{itemize}
   \item{Legg vekt på refleksjon, analytisk og kritisk tenkning}
   \item{Pek ut de viktigste programmeringsferdighetene du har utviklet}
   \item{Hvordan har programmering styrket forståelsen av økonomiske eller andre faglige aspekter?}
   \item{Har noen av oppgavene eller casene bidratt forbedret forståelse av økonomiske problemstillinger?}
   \item{Hvilke utfordringer møtte du, og hvordan håndterte du disse?}
   \item{Har du lært noe som er relevant for en karriere du ser for deg?}
   \item{Du kan diskutere potensielle anvendelser}
   \item{Identifiser noen ting som var vellykkede, og hvorfor}
   \item{Hvordan kan du bygge på det du har lært i fremtiden?}
   \item{Er det kanskje noe i oppgavene du har levert du kunne ønske å forbedre?}
   \item{Diskuter og grei ut om hva du har lært}
   \item{Du kan reflektere over egen tilnærming til å lære programmering}
   \item{Det vil være lurt å knytte refleksjoner til konkrete eksempler i mappen}
   \item{Refleksjonsnotat er en personlig tekst, så vær ærlig og åpen}
   \item{Det kan være lurt med en oppsummering og en slags konklusjon på slutten}
\end{itemize}

Du trenger naturligvis ikke å oppfylle alle punktene over, men de kan være til hjelp for å komme i gang.

~\\
\textit{Refleksjonsnotatet burde være på rundt 750-1500 ord}

\section*{Inspera}

Det åpnes for innlevering på Inspera mellom 18.12 og 22.12.
Hele mappen med kode og tekst pakkes sammen til en zip-fil og lastes opp i inspera.

Vi legger ut litt info og lenker på blackboard om hvordan du pakker sammen filer til et zip-arkiv, og hvor dere finner Jupyter-notebook filene deres


\end{document}
